\section{Policy Learning: método de Policy Gradient}

Para superar las limitaciones de Q-Learning, esta sección presenta otra clase de algoritmos de Deep RL—\textbf{Policy Gradient} (gradiente de política), que aprende directamente la función de política $\pi(s)$, sin necesidad de aprender la Q-function intermedia.

\subsection{Idea central de Policy Gradient}

\begin{figure}[H]
\centering
\includegraphics[width=0.85\textwidth]{slides-images/slide-38.jpg}
\caption{Política en un espacio de acciones discreto: salida de la distribución de probabilidad de cada acción $P(a|s)$\protect\footnotemark}
\end{figure}
\footnotetext{Slides, página 38.}

A diferencia de DQN, el método Policy Gradient ya no estima el valor Q de cada acción, sino que produce directamente una \textbf{distribución de probabilidad sobre las acciones}. La entrada de la red es el estado $s$, y la salida es $P(a|s)$—la probabilidad de tomar cada acción en ese estado.

Para un espacio de acciones \textbf{discreto}, la salida es una distribución de probabilidad softmax (similar a una tarea de clasificación):

$$P(a = \text{left} | s) = 0.6, \quad P(a = \text{stay} | s) = 0.25, \quad P(a = \text{right} | s) = 0.15$$

\subsection{Manejo de espacios de acciones continuos}

La mayor ventaja de Policy Gradient es que puede manejar de forma natural \textbf{espacios de acciones continuos}.

\begin{figure}[H]
\centering
\includegraphics[width=0.95\textwidth]{slides-images/slide-40.jpg}
\caption{Policy Gradient en espacios de acciones continuos: la red produce los parámetros de una distribución gaussiana $(\mu, \sigma^2)$\protect\footnotemark}
\end{figure}
\footnotetext{Slides, página 40.}

Para un espacio de acciones continuo, la red ya no produce la probabilidad de cada acción, sino que produce los \textbf{parámetros de una distribución de probabilidad}. Por ejemplo, supongamos que la acción sigue una distribución gaussiana:

$$P(a|s) = \mathcal{N}(\mu, \sigma^2)$$

La red produce la media $\mu$ y la varianza $\sigma^2$, y luego se muestrea de esa distribución para obtener una acción concreta:

$$\pi(s) \sim P(a|s) = \mathcal{N}(\mu, \sigma^2)$$

\begin{importantbox}{Ventajas clave de Policy Gradient}
\begin{itemize}
    \item \textbf{Manejo de espacios de acciones continuos}: al producir los parámetros de una distribución de probabilidad (como $\mu$ y $\sigma^2$ de una distribución gaussiana), puede manejar de forma natural problemas de control continuo
    \item \textbf{Aprendizaje de una política estocástica}: la salida es una distribución de probabilidad en lugar de una acción determinista, lo que da soporte natural a Stochastic Policy
    \item \textbf{Optimización directa del objetivo}: optimiza directamente la política para maximizar el retorno esperado, sin necesidad de aprender la Q-function intermedia
\end{itemize}
\end{importantbox}

\subsection{Algoritmo de entrenamiento de Policy Gradient}

\begin{figure}[H]
\centering
\includegraphics[width=0.85\textwidth]{slides-images/slide-44.jpg}
\caption{Flujo de entrenamiento de Policy Gradient: algoritmo de entrenamiento de cinco pasos\protect\footnotemark}
\end{figure}
\footnotetext{Slides, página 44.}

El entrenamiento de Policy Gradient sigue el siguiente flujo de cinco pasos:

\begin{enumerate}
    \item \textbf{Inicializar el Agent}: inicializar de forma aleatoria los parámetros de la red de política
    \item \textbf{Ejecutar la política hasta la terminación}: hacer que el Agent ejecute un episode completo en el entorno siguiendo la política actual
    \item \textbf{Registrar todos los (state, action, reward)}: recopilar todos los datos de trayectoria de todo el episode
    \item \textbf{Reducir la probabilidad de las acciones de baja recompensa}: para las acciones que llevan a un retorno bajo, reducir la probabilidad de que sean elegidas
    \item \textbf{Aumentar la probabilidad de las acciones de alta recompensa}: para las acciones que llevan a un retorno alto, aumentar la probabilidad de que sean elegidas
\end{enumerate}

\subsection{Función de pérdida de Policy Gradient}

El núcleo matemático del entrenamiento consiste en definir una función de pérdida adecuada:

$$\text{loss} = -\log P(a_t | s_t) \cdot R_t$$

\begin{figure}[H]
\centering
\includegraphics[width=0.95\textwidth]{slides-images/slide-46.jpg}
\caption{Función de pérdida y fórmula de actualización del gradiente de Policy Gradient\protect\footnotemark}
\end{figure}
\footnotetext{Slides, página 46. Williams Reinforcement Learning 1992.}

\begin{itemize}
    \item $-\log P(a_t | s_t)$: el logaritmo negativo de la verosimilitud (log-likelihood) de elegir la acción $a_t$
    \item $R_t$: el retorno acumulado descontado a partir del instante $t$
\end{itemize}

La regla de actualización por descenso de gradiente es:

$$w' = w - \nabla \text{loss} = w + \nabla \log P(a_t | s_t) \cdot R_t$$

\begin{importantbox}{Comprensión intuitiva de la función de pérdida de Policy Gradient}
El diseño de esta función de pérdida es muy ingenioso:
\begin{itemize}
    \item Cuando $R_t > 0$ (retorno alto), el signo negativo del loss hace que la actualización del gradiente \textbf{aumente} $P(a_t|s_t)$, es decir, incremente la probabilidad de esa acción
    \item Cuando $R_t < 0$ (retorno bajo o penalización), la actualización del gradiente \textbf{reduce} $P(a_t|s_t)$, es decir, disminuye la probabilidad de esa acción
    \item El valor absoluto de $R_t$ determina la magnitud de la actualización—cuanto mayor es el retorno, más fuerte es el ajuste de la política
\end{itemize}
Esta es la idea central del algoritmo REINFORCE (Williams, 1992).
\end{importantbox}

\begin{warningbox}{El problema de alta varianza de Policy Gradient}
Un problema conocido del método Policy Gradient es la \textbf{alta varianza} (High Variance). Como se usa el muestreo de Monte Carlo para estimar el retorno, la fluctuación del retorno entre distintos episode puede ser muy grande, lo que genera un ruido considerable en la estimación del gradiente e inestabilidad en el entrenamiento. Los métodos habituales para mitigar esto incluyen: usar un baseline para reducir la varianza (como en los métodos Actor-Critic), aumentar el número de muestras, y usar reward normalization, entre otros.
\end{warningbox}

\subsection{Resumen de la sección}

El método Policy Gradient aprende directamente la función de política y optimiza la política mediante «aumentar la probabilidad de las acciones de alto retorno y reducir la probabilidad de las acciones de bajo retorno». Su mayor ventaja es que puede manejar espacios de acciones continuos y aprender políticas estocásticas. La función de pérdida se construye a partir del producto entre el logaritmo de la verosimilitud de la acción y el retorno descontado, y el retorno esperado se optimiza mediante ascenso de gradiente.

